\chapter{Fondamenti teorici}
\label{cap:fondamenti}

Questo capitolo introduce da zero le nozioni necessarie a leggere il resto della
relazione. Chi ha già familiarità con modelli linguistici autoregressivi,
adattamento a rango basso e gradiente di policy può passare direttamente al
Capitolo~\ref{cap:dataset}.

L'ordine è quello delle dipendenze logiche: prima come si rappresenta il testo
(tokenizzazione), poi che cosa è un modello linguistico e come si addestra
(entropia incrociata), poi come si adatta un modello grande con poche risorse
(LoRA e quantizzazione), poi come si genera testo (decodifica), infine come si
addestra con un segnale di qualità invece che con un riferimento
(reinforcement learning e gradiente di policy).

\section{Tokenizzazione: il testo come sequenza di simboli}

Un modello linguistico non lavora su caratteri né su parole, ma su \emph{token}:
elementi di un vocabolario finito $\vocab$ costruito una volta per tutte a
partire da un grande corpus. Un token può essere una parola intera, un
frammento di parola, un segno di punteggiatura o uno spazio.

Formalmente, il tokenizzatore è una funzione
\begin{equation}
  \mathrm{tok}: \Sigma^* \longrightarrow \vocab^*
\end{equation}
dove $\Sigma$ è l'alfabeto dei caratteri e $\vocab^*$ l'insieme delle sequenze
finite di token. Si richiede che $\mathrm{tok}$ sia invertibile
(\emph{lossless}): esiste $\mathrm{detok}$ tale che
$\mathrm{detok}(\mathrm{tok}(s)) = s$ per ogni stringa $s$.

Questo dettaglio, apparentemente pedante, è centrale in questo progetto. La
parola gloss \texttt{DESC-NOT} non è un token: viene spezzata in più token, per
esempio \texttt{DESC}, \texttt{-}, \texttt{NOT}. Ne segue che un vincolo
espresso a livello di \emph{parola} (``emetti solo gloss appartenenti al
vocabolario'') deve essere tradotto in un vincolo a livello di \emph{token}
(``emetti solo token che siano prefisso valido di almeno un gloss''). È
esattamente il problema che il Trie del Capitolo~\ref{cap:decoding} risolve.

Un secondo dettaglio: molti tokenizzatori moderni codificano lo spazio come
parte del token successivo. Il token che rappresenta \texttt{"HOUSE"} a inizio
frase e quello che rappresenta \texttt{"~HOUSE"} in mezzo a una frase sono
identificatori \emph{diversi}. Anche questo ha una conseguenza diretta
(la doppia radice del Trie, sezione~\ref{sec:trie-dual-root}).

\section{Modelli linguistici autoregressivi}

Sia $y = (y_1, \dots, y_T)$ una sequenza di token. Un modello linguistico
\emph{autoregressivo} fattorizza la probabilità congiunta della sequenza usando
la regola della catena:
\begin{equation}
  p_\theta(y) \;=\; \prod_{t=1}^{T} p_\theta(y_t \mid y_{<t}),
  \qquad y_{<t} = (y_1, \dots, y_{t-1}).
\end{equation}

Nel caso condizionale che ci interessa (dato un ingresso $x$, l'inglese,
produrre un'uscita $y$, il gloss) la fattorizzazione diventa
\begin{equation}
  \label{eq:ar-cond}
  p_\theta(y \mid x) \;=\; \prod_{t=1}^{T} p_\theta(y_t \mid x, y_{<t}).
\end{equation}

Il modello calcola, a ogni passo $t$, un vettore di \emph{logits}
$z_t \in \mathbb{R}^{|\vocab|}$, cioè un punteggio non normalizzato per ogni
token del vocabolario. La distribuzione di probabilità si ottiene applicando la
funzione softmax:
\begin{equation}
  \label{eq:softmax}
  p_\theta(y_t = v \mid x, y_{<t})
  \;=\;
  \frac{\exp(z_{t,v})}{\sum_{u \in \vocab} \exp(z_{t,u})}.
\end{equation}

Il denominatore, detto \emph{funzione di partizione}, garantisce che le
probabilità sommino a uno. Il fatto che la somma sia su \emph{tutto} il
vocabolario è ciò che rende non banale imporre un vincolo: se si vietano alcuni
token, il denominatore va ricalcolato sull'insieme ammesso.

\subsection{Architettura: Transformer, in breve}

L'architettura impiegata (Qwen2.5, si veda il Capitolo~\ref{cap:modello}) è un
Transformer \emph{decoder-only}. Non è necessario, per leggere questa relazione,
conoscerne i dettagli; bastano tre fatti.

\begin{enumerate}
  \item Ogni token viene trasformato in un vettore (\emph{embedding}) e passa
        attraverso una pila di blocchi identici nella struttura.
  \item Ogni blocco contiene un modulo di \emph{attenzione}, che mette in
        relazione la posizione corrente con tutte le posizioni precedenti, e un
        modulo \emph{feed-forward} (una rete densa a due strati con non
        linearità). I moduli dell'attenzione sono chiamati per convenzione
        \texttt{q\_proj}, \texttt{k\_proj}, \texttt{v\_proj}, \texttt{o\_proj};
        quelli feed-forward \texttt{gate\_proj}, \texttt{up\_proj},
        \texttt{down\_proj}. Sono esattamente i sette moduli su cui agisce
        l'adattamento LoRA descritto più avanti.
  \item L'attenzione è \emph{causale}: la posizione $t$ non può guardare le
        posizioni $> t$. È questa mascheratura che rende valida la
        fattorizzazione~\eqref{eq:ar-cond} e che permette di calcolare, in un
        unico passaggio, le probabilità di tutti i token di una sequenza nota.
\end{enumerate}

\section{Addestramento supervisionato: entropia incrociata}

Dato un insieme di esempi $\mathcal{D} = \{(x^{(i)}, y^{(i)})\}_{i=1}^N$,
l'obiettivo dell'addestramento supervisionato è massimizzare la
log-verosimiglianza dei riferimenti, ossia minimizzare la \emph{loss} di
entropia incrociata:
\begin{equation}
  \label{eq:ce}
  \mathcal{L}_{\text{CE}}(\theta)
  \;=\;
  -\frac{1}{N}\sum_{i=1}^{N} \sum_{t=1}^{T_i}
  \log p_\theta\!\left(y^{(i)}_t \,\middle|\, x^{(i)}, y^{(i)}_{<t}\right).
\end{equation}

Interpretazione: per ogni posizione si guarda quale probabilità il modello
assegna al token \emph{corretto} e si penalizza il logaritmo negativo di quella
probabilità. Se il modello assegna probabilità $1$ al token corretto il
contributo è $0$; se assegna probabilità prossima a $0$ il contributo diverge.

Due proprietà rilevanti nel seguito:

\begin{description}
  \item[Teacher forcing.] Nella~\eqref{eq:ce} il condizionamento è sul
        \emph{prefisso di riferimento} $y^{(i)}_{<t}$, non su quanto il modello
        avrebbe generato. Questo rende l'addestramento parallelizzabile su tutte
        le posizioni, ma crea una discrepanza fra addestramento e generazione
        (\emph{exposure bias}): in generazione il modello condiziona sui propri
        errori.
  \item[Mascheratura del prompt.] Nel nostro caso interessa solo che il modello
        impari a produrre il gloss, non a riprodurre l'inglese di ingresso. Le
        posizioni corrispondenti al prompt vengono quindi \emph{escluse} dalla
        somma. Nella pratica implementativa si assegna a quelle posizioni
        l'etichetta convenzionale $-100$, che le librerie usate interpretano come
        ``ignora'' (si veda il Capitolo~\ref{cap:sft}).
\end{description}

\subsection{Un obiettivo più debole: le etichette parziali}
\label{sec:partial-label}

Esiste una variante utile quando non si conosce il token corretto, ma si conosce
un \emph{insieme} $\allowed_t \subseteq \vocab$ che lo contiene. In quel caso si
può massimizzare la probabilità che il modello assegni all'insieme, invece che a
un singolo elemento:
\begin{equation}
  \label{eq:allowed-mass}
  \mathcal{L}_{\text{AM}}(\theta)
  \;=\;
  -\log \sum_{v \in \allowed_t} p_\theta(v \mid x, y_{<t}).
\end{equation}

Questa è la formulazione classica dell'apprendimento a etichette parziali
(\emph{partial-label learning}, \cite{cour2011partial}) e coincide con la
massimizzazione della verosimiglianza marginale (\emph{maximum marginal
likelihood}, \cite{guu2017mml}). Va sottolineata una proprietà che sarà decisiva
nel Capitolo~\ref{cap:ausiliari}: la~\eqref{eq:allowed-mass} è \emph{piatta}
sull'insieme ammesso. Qualunque ridistribuzione della massa \emph{interna} ad
$\allowed_t$ lascia la loss invariata. Ne segue che questo obiettivo può
migliorare la \emph{calibrazione} (quanta probabilità va sui token legali) ma non
può, da solo, migliorare l'exact match.

\section{Adattamento efficiente: LoRA e quantizzazione}

Aggiornare tutti i parametri di un modello linguistico richiede memoria per i
pesi, per i gradienti e per gli stati dell'ottimizzatore. Due tecniche riducono
drasticamente questo costo.

\subsection{LoRA: adattamento a rango basso}

L'idea di LoRA (\emph{Low-Rank Adaptation}, \cite{hu2021lora}) è che
l'aggiornamento necessario per adattare un modello pre-addestrato a un compito
specifico abbia \emph{rango intrinseco basso}. Invece di aggiornare una matrice
di pesi $W_0 \in \mathbb{R}^{d \times k}$, si congela $W_0$ e si apprende un
aggiornamento fattorizzato:
\begin{equation}
  \label{eq:lora}
  W \;=\; W_0 \;+\; \Delta W,
  \qquad
  \Delta W \;=\; \frac{\alpha}{r}\, B A,
\end{equation}
con $A \in \mathbb{R}^{r \times k}$, $B \in \mathbb{R}^{d \times r}$ e
$r \ll \min(d,k)$. Il numero di parametri addestrabili passa da $dk$ a
$r(d+k)$.

I due iperparametri hanno un significato preciso:
\begin{itemize}
  \item $r$ (\emph{rango}) fissa la capacità dell'adattamento;
  \item $\alpha$ (\emph{alpha}) è un fattore di scala; il rapporto $\alpha/r$
        nella~\eqref{eq:lora} serve a rendere l'ampiezza dell'aggiornamento
        indipendente dalla scelta di $r$, così che cambiare $r$ non richieda di
        ritarare il learning rate.
\end{itemize}
All'inizializzazione $B = 0$, quindi $\Delta W = 0$ e il modello adattato
coincide esattamente con quello di partenza.

\subsection{QLoRA: pesi quantizzati}

La quantizzazione riduce la precisione numerica con cui i pesi congelati sono
memorizzati. In QLoRA~\cite{dettmers2023qlora} i pesi base $W_0$ sono
memorizzati a 4 bit, mentre gli adattatori $A, B$ restano in precisione più
alta e sono gli unici a ricevere gradiente. Il risultato è che un modello che in
precisione a 16 bit non entrerebbe in memoria diventa addestrabile su una sola
GPU.

Il costo è un rumore di quantizzazione sui pesi base. Per il presente lavoro la
conseguenza operativa è che i risultati non sono bit-per-bit riproducibili
cambiando backend di quantizzazione, mentre lo sono a parità di backend e seed.

\section{Generazione: strategie di decodifica}

Data la distribuzione~\eqref{eq:softmax}, produrre una sequenza richiede una
regola di scelta. Le opzioni rilevanti qui sono tre.

\begin{description}
  \item[Greedy.] Si prende a ogni passo il token più probabile,
        $y_t = \arg\max_v p_\theta(v \mid x, y_{<t})$. È deterministico ma tende
        a produrre testo ripetitivo~\cite{holtzman2019,welleck2020unlikelihood}.
  \item[Campionamento con temperatura.] Si campiona da una distribuzione
        riscalata:
        \begin{equation}
          \label{eq:temperature}
          p^{(\tau)}_\theta(v) \;\propto\; \exp\!\left(\frac{z_{t,v}}{\tau}\right).
        \end{equation}
        Per $\tau \to 0$ si recupera il greedy; per $\tau = 1$ la distribuzione
        del modello; per $\tau > 1$ una distribuzione più piatta, quindi più
        varia. La temperatura è il principale strumento per controllare la
        diversità dei candidati, e nel reinforcement learning è ciò che
        determina se il gruppo di candidati contiene o no differenze su cui
        apprendere.
  \item[Decodifica vincolata.] Si azzera la probabilità dei token non ammessi
        prima della normalizzazione. Operativamente si pone
        $z_{t,v} \leftarrow -\infty$ per $v \notin \allowed_t$, cosicché
        \begin{equation}
          \label{eq:constrained}
          p_\theta(v \mid x, y_{<t}, \allowed_t)
          =
          \begin{cases}
            \dfrac{\exp(z_{t,v})}{\sum_{u \in \allowed_t}\exp(z_{t,u})}
              & v \in \allowed_t,\\[2ex]
            0 & \text{altrimenti.}
          \end{cases}
        \end{equation}
        Questa è la forma di generazione con \emph{lessico chiuso} adottata nel
        Capitolo~\ref{cap:decoding}, ed è imparentata con la decodifica
        grammaticalmente vincolata~\cite{geng2023gcd,park2024gad} e con la
        decodifica a vincoli lessicali~\cite{hokamp2017,post2018}.
\end{description}

Va osservato che la~\eqref{eq:constrained} \emph{non} è una semplice
riparametrizzazione: sposta massa di probabilità dai token vietati a quelli
ammessi in proporzione ai logits originali. Il modello non ``sa'' di essere
vincolato, dunque il vincolo può produrre sequenze che il modello non avrebbe
mai generato liberamente. È il motivo per cui il vincolo può \emph{peggiorare}
una metrica di sovrapposizione e contemporaneamente \emph{migliorare} la
validità strutturale.

\section{Strutture dati simboliche: il Trie}

Un \emph{Trie} (o albero dei prefissi) è un albero in cui ogni nodo rappresenta
un prefisso e ogni arco è etichettato con un simbolo. Dato un vocabolario finito
di stringhe $G$ (nel nostro caso: l'insieme dei gloss ammessi), il Trie
costruito su $G$ permette di rispondere in tempo proporzionale alla lunghezza
del prefisso alla domanda:
\begin{quote}
  ``dato il prefisso di token $s$, quali token possono seguire senza rendere
  impossibile il completamento in un elemento di $G$?''
\end{quote}

Formalmente, detto $\mathrm{Pref}(G)$ l'insieme dei prefissi di token degli
elementi di $G$, l'insieme ammesso è
\begin{equation}
  \label{eq:trie-allowed}
  \allowed(s) \;=\; \{\, v \in \vocab \;:\; s \cdot v \in \mathrm{Pref}(G) \,\}
\end{equation}
dove $\cdot$ denota la concatenazione. Il Trie è precisamente la struttura che
materializza $\mathrm{Pref}(G)$ in modo che la~\eqref{eq:trie-allowed} sia
calcolabile con un accesso a un dizionario.

Un Trie riconosce il linguaggio $G$, che è finito e quindi regolare. Un
automa a stack (\emph{pushdown automaton}) riconosce linguaggi context-free,
strettamente più espressivi. La scelta fra i due, nel
Capitolo~\ref{cap:decoding}, è motivata dal fatto che il linguaggio target di
questo progetto è \texttt{gloss*}, cioè una semplice chiusura di Kleene su un
insieme finito: un Trie è sufficiente, e un automa a stack sarebbe
sovradimensionato.

\section{Apprendimento per rinforzo: impostazione}

Nell'addestramento supervisionato si dispone del riferimento $y^{(i)}$. Nel
reinforcement learning si dispone invece di un \emph{segnale di qualità}, la
\emph{reward}, che valuta un'uscita prodotta dal modello stesso.

Si modella la generazione come un processo decisionale:
\begin{itemize}
  \item lo \emph{stato} al passo $t$ è il prefisso $(x, y_{<t})$;
  \item l'\emph{azione} è la scelta del token $y_t$;
  \item la \emph{policy} è il modello stesso, $\policy(y_t \mid x, y_{<t})$;
  \item la \emph{reward} $R(x, y) \in \mathbb{R}$ è assegnata alla sequenza
        completa (\emph{reward terminale}).
\end{itemize}

L'obiettivo è massimizzare la reward attesa:
\begin{equation}
  \label{eq:rl-objective}
  J(\theta) \;=\; \mathbb{E}_{x \sim \mathcal{D}}\;
                  \mathbb{E}_{y \sim \policy(\cdot \mid x)}
                  \big[ R(x, y) \big].
\end{equation}

\subsection{Il gradiente di policy}

Il problema tecnico è che la variabile rispetto a cui si deriva, $\theta$,
compare nella distribuzione su cui si prende l'aspettazione. La soluzione
classica è l'identità del \emph{likelihood ratio}: poiché
$\nabla_\theta \policy(y) = \policy(y)\, \nabla_\theta \log \policy(y)$, si
ottiene
\begin{equation}
  \label{eq:pg}
  \nabla_\theta J(\theta)
  \;=\;
  \mathbb{E}_{y \sim \policy}\big[ R(x,y)\, \nabla_\theta \log \policy(y \mid x) \big].
\end{equation}
Questo è il teorema del gradiente di policy, alla base di REINFORCE
\cite{williams1992reinforce}. Il significato è diretto: si aumenta la
log-probabilità delle sequenze con reward alta e si diminuisce quella delle
sequenze con reward bassa.

\subsection{Baseline e vantaggio}

La stima~\eqref{eq:pg} è corretta ma ha varianza elevata. Si dimostra che
sottrarre una qualunque quantità $b$ indipendente dall'azione lascia il
gradiente inalterato in aspettazione, poiché
$\mathbb{E}_{y \sim \policy}[\nabla_\theta \log \policy(y)] = 0$. Si definisce
allora il \emph{vantaggio}
\begin{equation}
  \label{eq:advantage}
  A(x,y) \;=\; R(x,y) - b(x)
\end{equation}
e si sostituisce $R$ con $A$ nella~\eqref{eq:pg}. La scelta di $b(x)$ è ciò che
distingue gli algoritmi fra loro: PPO~\cite{schulman2017ppo}, impiegato anche
nella messa a punto di modelli linguistici con feedback
umano~\cite{ouyang2022instructgpt}, apprende una rete
separata (il \emph{critic}) che stima $b(x)$; GRPO~\cite{shao2024grpo} usa
invece la media delle reward di un gruppo di candidati generati per lo stesso
$x$, eliminando il critic. Il dettaglio è nel Capitolo~\ref{cap:grpo}.

Da~\eqref{eq:advantage} segue un fatto che ricorrerà come chiave interpretativa
di tutta la parte sperimentale: se \emph{tutti} i candidati per un dato $x$
ricevono la stessa reward, allora $A = 0$ per ciascuno e il gradiente è nullo.
Quel prompt non contribuisce all'apprendimento. Chiameremo \emph{gruppo morto}
un gruppo in questa condizione, e \emph{supporto} la frazione di gruppi che
producono un gradiente non nullo.

\subsection{Vincolo di prossimità e regolarizzazione KL}

Aggiornamenti troppo grandi degradano la policy in modo irreversibile. Si
introduce quindi un termine che penalizza l'allontanamento da una policy di
riferimento $\reference$ (tipicamente il modello prima
dell'addestramento per rinforzo):
\begin{equation}
  \label{eq:kl}
  J_{\beta}(\theta)
  \;=\;
  \mathbb{E}\big[A(x,y)\big]
  \;-\;
  \beta \, D_{\mathrm{KL}}\!\left(\policy \,\|\, \reference\right),
\end{equation}
dove $D_{\mathrm{KL}}(p\|q) = \sum_v p(v) \log \frac{p(v)}{q(v)} \ge 0$ è la
divergenza di Kullback-Leibler, nulla se e solo se $p = q$. Il coefficiente
$\beta$ regola quanto la policy è libera di allontanarsi dal riferimento.

\subsection{Perché la reward va progettata con sospetto}

Ottimizzare una reward imperfetta produce comportamenti che massimizzano la
misura senza migliorare l'obiettivo reale. È il fenomeno noto come
\emph{reward hacking} o \emph{specification gaming}
\cite{amodei2016concrete,gao2022overoptimization}; un caso documentato e
pertinente è la deriva verso risposte più lunghe quando la reward correla con la
lunghezza~\cite{singhal2023length}. La teoria classica del \emph{reward shaping}
\cite{ng1999shaping} caratterizza quali trasformazioni della reward preservano
la politica ottima; nella pratica di questo progetto la contromisura adottata è
diversa e più diretta: ogni componente di reward viene sottoposta a test
avversari costruiti a mano (Capitolo~\ref{cap:reward}).

\section{Nozioni di teoria dell'informazione}

Due quantità servono nel seguito per misurare quanta incertezza esista in un
compito.

L'\emph{entropia} di una variabile discreta $Y$ misura l'incertezza media, in
bit:
\begin{equation}
  \label{eq:entropy}
  H(Y) \;=\; -\sum_{y} p(y) \log_2 p(y).
\end{equation}

L'\emph{entropia condizionale} misura l'incertezza residua su $Y$ una volta noto
$X$:
\begin{equation}
  \label{eq:cond-entropy}
  H(Y \mid X) \;=\; -\sum_{x,y} p(x,y) \log_2 p(y \mid x).
\end{equation}

La differenza $H(Y) - H(Y \mid X) = I(X;Y)$ è l'informazione mutua: quanta
incertezza su $Y$ viene rimossa dall'osservazione di $X$. Nel
Capitolo~\ref{cap:dataset} queste quantità vengono stimate empiricamente sul
corpus per rispondere a una domanda concreta: quanta incertezza resta sul gloss
una volta noto il token inglese corrispondente? La risposta ($0{,}856$ bit su
$8{,}578$) quantifica in che senso il compito sia quasi deterministico.

\section{Metriche di valutazione: definizioni}

Le metriche usate nel seguito vanno definite qui, perché la loro forma esatta
determina l'interpretazione dei risultati.

\begin{description}
  \item[Exact match (EM).] Frazione di uscite identiche al riferimento dopo
        normalizzazione. È una metrica binaria e severa, ma non manipolabile.
  \item[ROUGE-L.] Basata sulla più lunga sottosequenza comune
        (\emph{longest common subsequence}, LCS) fra generato $g$ e riferimento
        $r$~\cite{lin2004rouge}. Con
        $P = \mathrm{LCS}(g,r)/|g|$ e $R = \mathrm{LCS}(g,r)/|r|$, la misura
        $F$ è
        \begin{equation}
          \label{eq:rouge}
          F_1 \;=\; \frac{2PR}{P+R}.
        \end{equation}
        Essendo una sottosequenza (non una sottostringa), l'ordine conta ma non
        la contiguità.
  \item[BLEU.] Media geometrica delle precisioni su $n$-grammi, con una
        penalità per uscite troppo brevi~\cite{papineni2002bleu}. Va distinta la
        variante \emph{a frase} (mediata sulle frasi) da quella \emph{a corpus}
        (con conteggi aggregati): danno numeri diversi e vanno riportate
        separatamente.
  \item[chrF.] Analoga a BLEU ma su $n$-grammi di \emph{caratteri}
        \cite{popovic2015chrf}; più robusta alla morfologia.
  \item[Pass@$k$.] Frazione di prompt per cui almeno una fra $k$ generazioni
        supera una soglia di correttezza. Detti $g^{(i)}_1, \dots, g^{(i)}_N$ le
        $N$ generazioni per il prompt $i$, $r^{(i)}$ il riferimento, $\theta$ la
        soglia e $\mathbb{1}[\cdot]$ la funzione indicatrice, la definizione
        adottata in questo lavoro è
        \begin{equation}
          \label{eq:passk}
          \mathrm{pass@}k
          \;=\;
          \frac{1}{M}\sum_{i=1}^{M}
          \mathbb{1}\!\left[\;\exists\, j \le k :\;
            \mathrm{ROUGE\text{-}L}\!\left(g^{(i)}_j, r^{(i)}\right) \ge \theta
          \right],
        \end{equation}
        dove $M$ è il numero di prompt. Si tratta di una \emph{media empirica}
        sulle prime $k$ generazioni di ciascun prompt, con $\theta = 0{,}3$.

        Va segnalato, perché è fonte di confusione, che questa \textbf{non} è la
        forma combinatoria $1 - \binom{n-c}{k} / \binom{n}{k}$ diffusa in
        letteratura, la quale stima senza distorsione la probabilità di successo
        su $k$ estrazioni a partire da $n > k$ campioni osservati. Quella forma
        non è impiegata qui. La differenza è discussa nel
        Capitolo~\ref{cap:metriche}, insieme a un secondo stimatore empirico
        della stessa quantità che il codice riporta accanto al primo.
\end{description}

Un avvertimento che il Capitolo~\ref{cap:metriche} svilupperà in dettaglio:
tutte le metriche di questa lista, tranne l'exact match, sono metriche di
\emph{sovrapposizione}. Su un compito in cui l'uscita corretta è in larga parte
una copia dell'ingresso, esse assegnano punteggi alti a sistemi che non hanno
imparato nulla. Per questo il progetto introduce una metrica complementare, la
\emph{non-copy-token accuracy}, che valuta solo le posizioni in cui il gloss
\emph{differisce} dal testo di partenza.

\section{Riepilogo della notazione}

\begin{center}
\begin{tabular}{ll}
\toprule
Simbolo & Significato \\
\midrule
$x$ & frase inglese di ingresso (prompt) \\
$y$ & sequenza di gloss generata \\
$\vocab$ & vocabolario di token del modello \\
$\allowed_t$ & insieme dei token ammessi al passo $t$ \\
$\policy$ & policy corrente (il modello addestrato) \\
$\reference$ & policy di riferimento (modello prima del RL) \\
$z_t$ & logits al passo $t$ \\
$R(x,y)$ & reward della sequenza completa \\
$A(x,y)$ & vantaggio \\
$G$ & numero di candidati per gruppo (GRPO) \\
$r, \alpha$ & rango e scala di LoRA \\
$\beta$ & coefficiente della penalità KL \\
$\tau$ & temperatura di campionamento \\
\bottomrule
\end{tabular}
\end{center}
